\section{Geometric and physical setup}
\subsection{Bondi frame and hard gyroscopic memory}
Let $(u,r,\theta^A)$ be Bondi coordinates, with
\[
g_{AB}=r^2\gamma_{AB}+rC_{AB}+O(1),\qquad N_{AB}=\partial_u C_{AB}.
\]
As the reference for the hard/soft split, we fix the standard initially shear-free vacuum frame
\begin{equation}
C_{AB}(u=-\infty)=0.
\label{eq:early-shear-free}
\end{equation}
The final shear $\Delta C_{AB}$ may be nonzero, so ordinary displacement memory is not excluded.

Following the conventions of Seraj--Oblak and Faye--Seraj, define $\widetilde C_{AB}=\epsilon_{CA}C^C{}_B$. The hard part of the gyroscope precession is
\begin{equation}
\widehat\Omega_{\rm H}=-\frac18N_{AB}\widetilde C^{AB},
\label{eq:hard-precession}
\end{equation}
and the physical orientation angle in the radiation zone is
\begin{equation}
\Delta\Psi_{\rm H}=\frac{1}{r^2}\bar\Phi_{\rm H}+O(r^{-3}),\qquad
\bar\Phi_{\rm H}=-\frac18\int du\,N_{AB}\widetilde C^{AB}.
\label{eq:hard-memory}
\end{equation}
We refer to $\Delta\Psi_{\rm H}$ as the hard/helicity gyroscopic gravitational memory. It is neither the full gyroscopic memory nor the spin memory \cite{SerajOblak2023,FayeSeraj2025}.

\subsection{Null Jacobi system}
Consider an auxiliary null-geodesic segment $\gamma:[\lambda_0,\lambda_1]\to M$ crossing the radiation in a four-dimensional Lorentzian spacetime. Let $\lambda$ be an affine parameter, $k=d\gamma/d\lambda$, and $L=\lambda_1-\lambda_0>0$. The segment $\gamma$ is not the timelike gyroscope worldline. We also exclude an outgoing null generator comoving with the wave on $u={\rm const.}$; instead we choose a radiation-crossing null direction with $du/d\lambda\neq0$ so that the leading radiative curvature is sampled. We further work in a local radiation-zone patch with $L$ fixed and $L/r\to0$ as $r\to\infty$, which pushes angular and radial drift effects along $\gamma$ beyond leading order in $1/r$. The screen basis is adapted so that, in this local patch, its transverse components agree with the Bondi/TT polarization dyad. These are geometric hypotheses of the main theorem, not a direction-independent statement for arbitrary null directions at finite $r$.

At leading asymptotically flat, locally plane-wave order, $du/d\lambda$ is constant. Using positive affine freedom we set $du/d\lambda=1$, so that the retarded-time width of the selected stationary-to-stationary window equals $L$. Let $e_A$ be a parallel-propagated orthonormal screen basis and define
\begin{equation}
T_{AB}(\lambda)=-R(e_A,k,e_B,k).
\label{eq:tidal}
\end{equation}
If $\Ric(k,k)=0$, then $T$ is a symmetric trace-free $2\times2$ matrix.

With $t=(\lambda-\lambda_0)/L$, define the dimensionless profile
\[
A(t)=L^2T(\lambda_0+Lt).
\]
The Jacobi fundamental matrix obeys
\begin{equation}
\frac{d\Phi}{dt}=\begin{pmatrix}0&I_2\\ A(t)&0\end{pmatrix}\Phi,
\qquad \Phi(0)=I_4.
\label{eq:jacobi}
\end{equation}
Writing the free-propagation matrix as $F_0(s)=\begin{psmallmatrix}I_2&sI_2\\0&I_2\end{psmallmatrix}$, define the scattering matrix with the flat drift removed by
\begin{equation}
S_\gamma=F_0(-1)\Phi(1).
\label{eq:S}
\end{equation}
For the block decomposition $S_\gamma=(S_{ij})$, define the real scalar $\omega_\gamma$ through
\begin{equation}
\operatorname{skew}(S_{11}+S_{22})=\omega_\gamma J_2,
\qquad J_2=\begin{pmatrix}0&-1\\1&0\end{pmatrix}.
\label{eq:omega-def}
\end{equation}
We call $\omega_\gamma$ the nonlinear Jacobi memory/scattering channel. In the plane-wave memory-tensor language, it corresponds to the antisymmetric combination of the diagonal Jacobi blocks \cite{Flanagan2020,Harte2025}.

\subsection{Bel--Robinson curvature action}
We use metric signature $(-+++)$ and normalize the Bel--Robinson tensor as
\begin{equation}
B_{abcd}=\frac14\left(C_{aecf}C_b{}^e{}_d{}^f+{}^\star C_{aecf}\,{}^\star C_b{}^e{}_d{}^f\right).
\label{eq:BR-def}
\end{equation}
On a vacuum null screen,
\begin{equation}
B(k,k,k,k)=|\Psi_0|^2=\frac12\tr(T^2).
\label{eq:BR-tidal}
\end{equation}
We therefore define
\begin{equation}
Q_\gamma:=L^3\int_{\lambda_0}^{\lambda_1}B(k,k,k,k)\,d\lambda
=\frac12\int_0^1\tr(A(t)^2)\,dt.
\label{eq:Q}
\end{equation}
Under a positive affine reparametrization $\lambda'=a\lambda+b$ with $a>0$, one has $L'=aL$, $k'=a^{-1}k$, $d\lambda'=a\,d\lambda$, and $B(k',k',k',k')=a^{-4}B(k,k,k,k)$. Hence $L'^3d\lambda'\,B(k'^4)=L^3d\lambda\,B(k^4)$. Thus the normalization $du/d\lambda=1$ merely selects a representative of the affine-equivalence class of the fixed radiation-crossing segment and does not spoil the affine invariance of $Q_\gamma$. Changing the segment endpoints, the radiation-crossing null direction, or the screen identification instead defines a different auxiliary-probe problem; $Q_\gamma$ is not window independent or direction independent.

\begin{remark}
$Q_\gamma$ is a dimensionless windowed null action quadratic in the Weyl curvature. It is not an exposure measured on the gyroscope worldline, the Bondi mass loss, Isaacson effective stress-energy, radiated energy, or detector energy.
\end{remark}

\section{Jacobi memory and the minimum-action problem}
Using the STF basis
\[
\sigma_3=\begin{pmatrix}1&0\\0&-1\end{pmatrix},\qquad
\sigma_1=\begin{pmatrix}0&1\\1&0\end{pmatrix},
\]
write
\begin{equation}
A(t)=\alpha(t)\sigma_3+\beta(t)\sigma_1,
\qquad Q_\gamma=\|\alpha\|_2^2+\|\beta\|_2^2.
\label{eq:control}
\end{equation}

\begin{definition}
Define
\begin{equation}
(K_0f)(t)=\frac12\int_0^1(t-s)|t-s|f(s)\,ds,
\qquad s_0:=\|K_0\|.
\label{eq:K0}
\end{equation}
The kernel is antisymmetric, so $K_0^*=-K_0$; it is also Hilbert--Schmidt and hence compact.
\end{definition}

\begin{lemma}[Small-amplitude expansion of the exact endpoint map]
\label{lem:endpoint-expansion}
For sufficiently small $A$, there exists $C>0$ such that
\begin{equation}
\omega_\gamma[A]=2\langle\alpha,K_0\beta\rangle+R_\omega[A],
\qquad |R_\omega[A]|\le C Q_\gamma[A]^2.
\label{eq:endpoint-expansion}
\end{equation}
\end{lemma}
\begin{proof}
Write each Jacobi column of Eq.~\eqref{eq:jacobi} as $X''=AX$. Then
\[
X(t)=X(0)+tX'(0)+(\mathcal V_A X)(t),\qquad
(\mathcal V_A X)(t):=\int_0^t(t-s)A(s)X(s)\,ds.
\]
Since the interval has unit length,
\[
\|\mathcal V_A X\|_\infty
\le \|A\|_{L^1,F}\|X\|_\infty
\le \|A\|_{L^2,F}\|X\|_\infty.
\]
Hence, for sufficiently small $\|A\|_{L^2,F}$, the Neumann--Volterra series converges absolutely and the terminal values $X(1)$ and $X'(1)$, and therefore $\Phi(1)$ and $\omega_\gamma[A]$, are analytic in $A$ near the origin. Let
\[
\mathcal V=\operatorname{span}\{\sigma_1,\sigma_3\},\qquad
\mathcal E=\operatorname{span}\{I_2,J_2\}.
\]
Then $\mathcal V\mathcal V\subset\mathcal E$ and $\mathcal E\mathcal V\subset\mathcal V$. Consequently, any product of an odd number of $A(t_i)\in\mathcal V$ lies again in the symmetric space $\mathcal V$ and cannot contribute to the antisymmetric channel. The linear term is also symmetric. In the second-order time-ordered term, only the noncommuting part of the two STF basis matrices survives. Using $[\sigma_3,\sigma_1]=-2J_2$ yields the kernel $\frac12(t-s)|t-s|$ and therefore the leading term $2\langle\alpha,K_0\beta\rangle$. The next possible contribution is fourth order, so in a fixed neighborhood
\[
|R_\omega[A]|\le C_1\|A\|_{L^1,F}^4
\le C_1\|A\|_{L^2,F}^4
=C Q_\gamma^2.
\]
\end{proof}

\section{Sharp minimum-action theorem}
Define
\begin{equation}
\kappa_\Omega(\omega):=\inf\{Q_\gamma[A]:\omega_\gamma[A]=\omega\}.
\label{eq:value}
\end{equation}

\begin{theorem}[Sharp law for the unrestricted nonlinear Jacobi channel]
\label{thm:jacobi-sharp}
\begin{equation}
\boxed{\kappa_\Omega(\omega)=\frac{|\omega|}{\|K_0\|}+O(\omega^2),\qquad \omega\to0.}
\label{eq:sharp-jacobi}
\end{equation}
The coefficient is sharp for the exact endpoint problem.
\end{theorem}
\begin{proof}
By Cauchy--Schwarz and the arithmetic--geometric mean inequality,
\begin{equation}
2|\langle\alpha,K_0\beta\rangle|
\le s_0(\|\alpha\|_2^2+\|\beta\|_2^2)=s_0Q_\gamma.
\label{eq:quad-lower}
\end{equation}
Compactness implies that the norm of $K_0$ is attained. Choose a real unit vector $\beta_*$ in the top eigenspace of $-K_0^2$,
\[
-K_0^2\beta_*=s_0^2\beta_*,\qquad
\alpha_*:=K_0\beta_*/s_0.
\]
Then $\|\alpha_*\|=\|\beta_*\|=1$ and $\langle\alpha_*,K_0\beta_*\rangle=s_0$. For $A_*=\alpha_*\sigma_3+\beta_*\sigma_1$, one has $Q_\gamma[A_*]=2$ and a quadratic memory coefficient $2s_0$, saturating Eq.~\eqref{eq:quad-lower}.

For exact recovery, set $A_\rho=\rho A_*$. The preceding lemma gives
\begin{equation}
\omega_\gamma[A_\rho]=2s_0\rho^2+O(\rho^4),\qquad Q_\gamma[A_\rho]=2\rho^2.
\label{eq:recovery-ray}
\end{equation}
Analyticity then yields
\[
\frac{d}{d\rho}\omega_\gamma[A_\rho]=4s_0\rho+O(\rho^3)>0
\]
for sufficiently small $\rho>0$. Thus we do not invoke the usual inverse-function theorem at $\rho=0$, where the derivative vanishes; instead, strict one-sided monotonicity gives a unique $\rho(\omega)>0$ for every sufficiently small positive target. Hence
\[
Q_\gamma[A_{\rho(\omega)}]=\frac{\omega}{s_0}+O(\omega^2).
\]
For negative targets replace $\beta_*$ by $-\beta_*$; the same one-sided argument gives the negative branch.

Conversely, the recovery construction already implies $\kappa_\Omega(\omega)=O(|\omega|)$, so a minimizing sequence may be chosen with $Q_\gamma=O(|\omega|)$. Applying Eqs.~\eqref{eq:endpoint-expansion} and \eqref{eq:quad-lower} along that sequence gives
\[
|\omega|\le s_0Q_\gamma+C Q_\gamma^2.
\]
Since $Q_\gamma=O(|\omega|)$, this yields $Q_\gamma\ge |\omega|/s_0-O(\omega^2)$. Combining the lower and upper bounds proves Eq.~\eqref{eq:sharp-jacobi}.
\end{proof}

A Gauss--Legendre Nystr\"om calculation gives
\begin{equation}
\|K_0\|=0.09103843824\ldots,\qquad \|K_0\|^{-1}=10.98437121\ldots.
\label{eq:numeric-k0}
\end{equation}
These decimal values are numerical validation only and do not define the theorem constant.